\section{Finite certificate selection before estimation}
\label{sec:finite-order-policy}

I now execute the certificate as a stopping procedure, rather than computing
all retained coefficients before selecting a price. This section distinguishes
three decisions: a certified stopping test, a sufficient screen which avoids
unnecessary scalar work, and a prescribed degree limit. None requires
convergence of the infinite Taylor series. The finite coefficient and
remainder identities remain those in
Theorem~\ref{thm:finite-gaussian-price-certificate}.

\subsection{Outward finite-order intervals and stopping}

For a candidate degree $k$, let $[a_k,b_k]$ be the outward arithmetic
enclosure of $\sum_{j=0}^k c_j$, and let $\overline R_k\geq R_k$ be a
nonnegative rational upper bound on its proved expected remainder. I form
\begin{equation}
 I_k=[\max\{0,a_k-\overline R_k\},
          \min\{H,b_k+\overline R_k\}].
 \label{eq:finite-policy-interval}
\end{equation}
The target belongs to both the Taylor enclosure and $[0,H]$, so their
intersection is nonempty and contains it. If the exact rational half-width
of $I_k$ is at most a requested undiscounted tolerance $\tau$, the midpoint
prices with deterministic error at most $\tau$. Testing this condition is
valid at every finite degree, independently of the success of other degrees.
There is no optional-stopping probability argument: each candidate is a
deterministic price certificate for the same target.

The full-search policy evaluates coefficients through degree eight once,
forms all nine partial intervals and selects the narrowest. The screened
policy computes $\overline R_k$ first. For $k<8$, it evaluates the
coefficient sum only when $\overline R_k\leq\tau$. It then applies the
actual half-width test, including coefficient arithmetic. At degree eight
it evaluates the coefficient sum even if the screen fails, so a valid
unresolved price interval is retained. A qualifying remainder does not
guarantee acceptance when arithmetic width consumes the remaining slack.
Conversely, clipping to $[0,H]$ may allow an interval to certify even with
$\overline R_k>\tau$. The screen is therefore sufficient for spending
coefficient work, not a proof of the globally first certifying degree.

\begin{proposition}[Certificate-preserving finite selection]
\label{prop:finite-certificate-selection}
Under the hypotheses of Theorem~\ref{thm:finite-gaussian-price-certificate},
each accepted finite-order midpoint has deterministic error at most $\tau$.
If no order is accepted, the order-limit interval remains a valid price
enclosure. Replacing full search by the specified remainder screen changes
the executed work and possibly the selected enclosure, but does not change
the validity of any returned certificate.
\end{proposition}
\begin{proof}
The coefficient enclosure and remainder theorem imply membership of the
target in $[a_k-\overline R_k,b_k+\overline R_k]$. Payoff boundedness gives
membership in $[0,H]$. Equation~\eqref{eq:finite-policy-interval} is exactly
their intersection. For any real number in an interval, its distance from
that interval's midpoint is no larger than its half-width. The acceptance
condition proves the asserted error. Skipping a coefficient calculation
at another degree asserts no interval for that degree and cannot alter
this implication. The explicitly evaluated terminal degree supplies the
unresolved enclosure when acceptance does not occur.
\end{proof}

Each moment degree $j$ has $\binom{n+j-1}{j}$ multinomial terms. Summing
these counts through degree $d$ gives $\binom{n+d}{d}$, by induction using
Pascal's identity. I retain a geometry bank containing each term's exact
multinomial weight, exposure and validated exponential. Extending the bank
for a failed screen incurs its full new term count. Later degrees reuse
the existing terms without charging them twice. A coefficient evaluation
through degree $k$ calls the Gaussian derivative kernel at all terms
through that degree, hence $\binom{n+k}{k}$ calls in the implemented
routine. If several degrees require coefficient evaluation, every such
call is charged; the screening algorithm does not silently borrow
unexecuted full-search derivatives. Geometry sharing across caps and
common loadings is allowed only when the exact centred geometry matches.
Cold-contract experiments own separate banks.

\subsection{Why the raw remainder must remain finite-order}

\begin{proposition}[Eventual growth of this Gaussian remainder bound]
\label{prop:raw-remainder-growth}
Suppose at least one centred loading $s_i$ is nonzero. For every fixed
$H>0$ and $\lambda>0$, the raw bound $R_k$ in
equation~\eqref{eq:taylor-price-remainder} tends to infinity as
$k\to\infty$.
\end{proposition}
\begin{proof}
Put $r=k+1$. The signed-Stirling recurrence gives
$s(r,1)=(-1)^{r-1}(r-1)!$ by induction from $s(1,1)=1$.
The $\ell=1$ term of the nonnegative sum defining $C_r$ therefore implies
$C_r(\lambda)\geq(r-1)!\phi(0)/\lambda$. Choose any $a>1$ and an index
with $s_i\ne0$. Since $s_i\cdot Y$ is a normal variable of positive
variance, its upper tail has positive probability at every finite threshold.
In particular, the event
\[
 s_i\cdot Y\geq \log((a+1)/w_i)-d_i
\]
has probability $p_a>0$ and implies
$T\geq w_i e^{d_i+s_i\cdot Y}\geq a+1$. Since $T-1\geq0$,
$\mathbb E[(T-1)^r]\geq p_a a^r$. Combining these bounds gives
\[
 R_{r-1}\geq {H\phi(0)p_a\over\lambda r}a^r
       \longrightarrow\infty.
\]
The last limit follows because $a>1$: the ratio of consecutive terms
$a^r/r$ tends to $a$, which exceeds one.
\end{proof}

This is a derived limitation of the particular uniform absolute-derivative
remainder used here. It is not a proof that the actual finite approximation
errors diverge, that another signed remainder cannot improve, or that no
other classical certificate works. The price intervals clipped to $[0,H]$
remain bounded. The proposition strengthens the reason for retaining a
finite degree search: ever higher degree cannot be justified by assuming
that this raw bound eventually vanishes. The visualised early improvements
and later failures are compatible with this statement. I do not claim
priority for the elementary tail argument or for generic lognormal moment
growth; its role is to characterise this explicitly defined certificate.

\subsection{A monotone geometric-control certificate before Taylor screening}
\label{sec:monotone-before-taylor}

I place the existing monotonicity and first-moment bound before the finite
Taylor remainder screen. This changes the order in which I pay for certified
computations. It does not assert a new Jensen inequality, a new geometric
control or convergence of an infinite Taylor series.

Let $w_i>0$, $\sum_{i=1}^n w_i=1$, and let $Y\sim N(0,I_p)$ and
$Z\sim N(0,1)$ be independent. For fixed real $\mu_i$ and $R_i\in\mathbb R^p$,
positive $\lambda$, spot $S_*=100$ and strike $K=100$, write
\[
 S_i=S_*\exp(\mu_i+R_i\cdot Y+\lambda Z),\quad
 A=\sum_iw_iS_i,\quad G=\prod_iS_i^{w_i}.
\]
Put $\mu_G=\sum_iw_i\mu_i$, $R_G=\sum_iw_iR_i$,
$d_i=\mu_i-\mu_G$ and $s_i=R_i-R_G$. Then
\[
 T(Y)={A\over G}=\sum_iw_i\exp(d_i+s_i\cdot Y),
 \qquad \sum_iw_i(d_i+s_i\cdot Y)=0.
\]
Weighted arithmetic-geometric mean therefore gives $T(Y)\geq1$ for every
$Y$. Its Gaussian first moment is finite and is given exactly by
\begin{equation}
 M_1=\mathbb E[T]=\sum_iw_i\exp\{d_i+\|s_i\|^2/2\}\geq1.
 \label{eq:hierarchy-ratio-first-moment}
\end{equation}
The inequality also follows by taking expectations of the pointwise bound.

For cap $H>0$, define $h_H(x)=\min\{(x-K)^+,H\}$ and the conditional scalar
kernel
\[
 g_{H,\lambda}(m)=\mathbb E_Z[h_H(S_*e^{m+\lambda Z})].
\]
For every nonnegative $x$,
$h_H(x)=\int_K^{K+H}\mathbf1_{\{x>k\}}\,dk$; the choice of strict endpoint
does not affect this integral. Tonelli's theorem gives
\[
 g_{H,\lambda}(m)=\int_K^{K+H}
 \Phi\left({m-\log(k/S_*)\over\lambda}\right)\,dk.
\]
The integration interval is finite and the derivative integrand is bounded
by $\phi(0)/\lambda$. Differentiation under this integral is therefore
justified by domination, and
\[
 0\leq g'_{H,\lambda}(m)
 ={1\over\lambda}\int_K^{K+H}
 \phi\left({m-\log(k/S_*)\over\lambda}\right)\,dk
 \leq {H\phi(0)\over\lambda}.
\]
A second bound follows directly from the payoff. It is Lipschitz in its
price argument, and its derivative with respect to the log-price exists
away from the two strikes and equals
$x\mathbf1_{\{K<x<K+H\}}$. This quantity is at most $K+H$. For a bounded
interval of $m$, the difference quotients are dominated by a finite
constant times $e^{\lambda Z}$, whose Gaussian expectation is finite.
The strike points have probability zero under the continuous normal law.
Dominated convergence consequently gives
$g'_{H,\lambda}(m)=\mathbb E[S_*e^{m+\lambda Z}
\mathbf1_{\{K<S_*e^{m+\lambda Z}<K+H\}}]\leq K+H$.
Thus $g_{H,\lambda}$ is increasing and globally $H L$-Lipschitz, where
\begin{equation}
 L=\min\left\{{\phi(0)\over\lambda},{K+H\over H}\right\}.
 \label{eq:hierarchy-log-lipschitz}
\end{equation}

Let $P_A=\mathbb E[h_H(A)]$ and $P_G=\mathbb E[h_H(G)]$ denote undiscounted
prices. Conditional on $Y$, their scalar log-price arguments differ by
$\log T(Y)\geq0$. The fundamental theorem of calculus and
\eqref{eq:hierarchy-log-lipschitz} imply
\[
 0\leq g_{H,\lambda}(\mu_G+R_G\cdot Y+\log T)
       -g_{H,\lambda}(\mu_G+R_G\cdot Y)
 \leq H L\log T.
\]
Since $0\leq\log T\leq T-1$ and $M_1<\infty$, this logarithm is integrable.
The conditional prices are themselves bounded by $H$. Taking expectations
and applying concavity of the logarithm gives the deterministic enclosure
\begin{equation}
 0\leq P_A-P_G\leq H L\,\mathbb E[\log T]
                    \leq H L\log M_1=:D_0.
 \label{eq:hierarchy-one-sided-gap}
\end{equation}
This proof requires positivity, positive weights and an increasing
log-Lipschitz conditional payoff kernel. If another outer law has a finite
first ratio moment, the inequality through $\log\mathbb E[T]$ remains
valid; its moment formula and its geometric-control computation must then
be certified for that law. The current implementation uses the stipulated
Gaussian law and does not transfer its Gaussian moment formula to another
model without that work.

I compute an outward upper enclosure $\overline L\geq L$, using the minimum
of the certified upper endpoint for $\phi(0)/\lambda$ and the exact rational
$(K+H)/H$. The validated degree-one moment ball contains $M_1$ and has a
strictly positive lower endpoint. Its logarithm therefore contains
$\log M_1$. The upper endpoint of
$H\overline L\log(M_1\text{ ball})$, clamped below by zero, supplies a
rational $\overline D_0\geq D_0$. Neither a negative midpoint nor a
floating-point logarithm is substituted for this upper certificate.

Let $\tau=0.01$ be the undiscounted price-radius target. If
$\overline D_0>2\tau$, I do not evaluate the scalar geometric control at
this preliminary stage. The one-sided moment certificate alone has not
passed its sufficient width screen, so I continue to the prescribed finite
Taylor screening policy and charge the moment work already performed.
When $\overline D_0\leq2\tau$, I evaluate $P_G$ once. Here
$\mu_G+R_G\cdot Y+\lambda Z$ is normal with mean $\mu_G$ and variance
$\sigma_G^2=\|R_G\|^2+\lambda^2$. The validated scalar kernel at that
combined scale gives an interval $[g_-,g_+]$ containing $P_G$. I form
\begin{equation}
 I_0=\left[\max\{0,g_-\},
          \min\{H,g_++\overline D_0\}\right].
 \label{eq:hierarchy-one-sided-price-interval}
\end{equation}
Its lower endpoint is no larger than $P_A$, because $P_A\geq P_G$ and
$P_A\geq0$. Its upper endpoint is no smaller than $P_A$, because
$P_A\leq P_G+\overline D_0$ and $P_A\leq H$. Hence the interval is
nonempty and contains the target price.

I stop only when the exact rational half-width of $I_0$ is at most $\tau$.
The preliminary inequality $\overline D_0\leq2\tau$ does not alone imply
this final condition: the scalar-control interval has its own positive
width, and at the boundary that width can prevent certification. If the
actual half-width is too large, I retain the failed preliminary interval,
charge its one scalar call and proceed to the remainder screen. This
distinction is tested using an intentionally widened but valid control
enclosure. The midpoint of an accepted $I_0$ has deterministic price error
at most $\tau$. Discounting is applied afterwards under the existing
pricing convention; no statistical confidence allowance is consumed by
this deterministic certificate.

Each policy owns a fresh geometry bank. Moment degree $j$ has
$\binom{n+j-1}{j}$ multinomial terms, and extending a bank through degree
$d$ requires exactly $\binom{n+d}{d}$ exponential evaluations, including
the constant degree-zero term. The hierarchy first extends its bank through
degree one. A rejected preliminary screen does not discard that bank: the
Taylor fallback extends the same bank as needed, so previously computed
geometry is charged once. Coefficients are evaluated only when their
remainder permits the prescribed screen, or at the fixed degree limit to
retain an unresolved interval. A successful preliminary certificate costs
one scalar kernel call; an unsuccessful numerical control enclosure costs
that call in addition to every subsequent Taylor coefficient evaluation.

The retained comparison supplies the same 27 contracts, degree limit eight,
224-bit ball arithmetic, coefficient identities and price target to all
three policies. It compares finite executed choices. It establishes neither
a globally cheapest screening order nor a complexity lower bound for the
pricing problem. The one-sided step is an application of the existing
control certificate; any priority claim for that inequality would require
the separate literature comparison already stipulated in this paper.
